% Generated by escalation/run_bench_script/make_figures_tex.py -- do not edit.
% The caption text lives in that chart's plot script, next to the numbers it
% describes; re-run the script to update this file.
\begin{figure}[p]
\centering
\begin{subfigure}{\linewidth}
  \centering
  \panelplot{\plotdir/manager_vs_single_call_across_model_scale_lcb-100_1_pass_128k_max_tokens_reasoning_off_light.png}
  \caption{Accuracy by model scale.}
\end{subfigure}\\[0.7ex]
\begin{subfigure}{\linewidth}
  \centering
  \panelplot{\plotdir/manager_vs_single_call_across_model_scale_lcb-100_1_pass_128k_max_tokens_reasoning_off_bars_light.png}
  \caption{The same numbers as paired bars.}
\end{subfigure}
\caption{\textbf{Reasoning off:} 128k $\times$ 1-pass control. \footnotesize\normalfont
Fill: light = single call (one call, no tools), dark = with manager. 128k max tokens, reasoning off (ESCALATION\_OR\_REASONING=none), 1 pass per condition. $\Delta$ = with manager $-$ single call, in percentage points. With one pass per condition there is no within-model repeat, so no per-model p is shown. Each block under the axis heads with the model's size, not its name - 9B and 35B are too close on the log axis of (a) for a name to fit - and the legend maps fill colour to name. The curves in (a) connect each condition's own 4 points (monotone cubic): they interpolate, they are not fits. Tukey groups: models sharing a letter have indistinguishable $\Delta$ (HSD, $\alpha$ 0.05); HSD ignores that all models saw the same 100 problems, so it is conservative. Pairwise Tukey p: 9B vs 35B 0.35 n.s.; 9B vs 428B 0.6 n.s.; 9B vs 2.8T $7.6\times10^{-7}$; 35B vs 428B 0.021; 35B vs 2.8T $5.3\times10^{-11}$; 428B vs 2.8T $2.5\times10^{-4}$.}
\end{figure}
